%! Author = WelJo
%! Date = 4/30/2026

%! Einfügung und Veränderung von "Ziel und Forschungsinteresse" zu Fragestellung - Inhaltliche Anforderung (Bao)

\section{Fragestellung}
\label{sec:fragestellung}

%% Beginning ===========================

Die Berechnung der Erreichbarkeit von Points of Interest (POIs) ist ein multidimensionales Problem.
Während klassische Ansätze wie der \textbf{Loss of Serviceability (LoS) Index} oder der \textbf{Link Performance Index (LPIR)} primär strukturelle oder verkehrsflussbezogene Aspekte des Straßennetzes bewerten, bleibt die funktionale Ebene der Ziele oft unberücksichtigt.
Ein POI ist jedoch in einer Krisensituation nur dann erreichbar, wenn er sowohl physisch zugänglich als auch funktional nutzbar ist.

Die aktuelle Forschungslücke liegt in der isolierten Betrachtung dieser Systeme.
Bestehende Verfahren zur Analyse der \textbf{Robustness of Accessibility (RoA)} konzentrieren sich auf die Widerstandsfähigkeit des Straßennetzes gegenüber physischen Störungen.
Gleichzeitig existieren Modelle zur Simulation von Kaskadeneffekten in Versorgungsnetzen, die den funktionalen Zusammenbruch von Einrichtungen über \textit{Lifetimes} und \textit{Repair Times} beschreiben.
Es fehlt jedoch eine methodische Verknüpfung dieser Welten: Ein POI kann laut RoA-Index perfekt erreichbar sein, während er laut Kaskadenmodell bereits aufgrund eines Stromausfalls funktionsunfähig ist.

Aus dieser Lücke leitet sich die zentrale Forschungsfrage ab:
\begin{quote}
\textit{Wie beeinflussen Kaskadeneffekte infolge von Ausfällen in Versorgungsnetzen die funktionale Erreichbarkeit (RoA) kritischer POIs im zeitlichen Verlauf einer Krise?}
\end{quote}

Im Rahmen des Projekts werden folgende Hypothesen und Herausforderungen untersucht:
\begin{itemize}
\item \textbf{H1:} Die Integration von Stromausfall-Kaskaden führt zu einer signifikant geringeren Resilienz, als eine rein straßenbasierte RoA-Analyse vermuten lässt\textsuperscript{}.
\item \textbf{H2:} Durch die Modellierung von \textit{Lifetimes} lassen sich kritische Zeitfenster identifizieren, in denen ein POI trotz blockierter Wege funktional bleibt oder trotz freier Wege ausfällt.
\end{itemize}  

Das Ziel dieser Forschungsarbeit ist es, eine methodische und technische Brücke zwischen der straßenbasierten RoA-Analyse und zeitabhängigen Kaskadenmodellen zu schlagen.
Hierzu wird die bestehende Implementierung des RoA-Verfahrens um die Komponente der Elektrizitätsversorgung erweitert.
Das Ergebnis ist ein kombiniertes Analysewerkzeug, das die Resilienz nicht mehr nur als Erreichbarkeit von Orten, sondern als Verfügbarkeit von Diensten definiert.
Dies bietet eine präzisere Entscheidungsgrundlage für die Alarm- und Einsatzplanung im Katastrophenschutz.

%% Ending ==============================